\sezione{1}{L'idea: dalla folla alla foresta}{Parte I · il concetto generale}

\begin{inbreve}
Un \hl{Random Forest} è un insieme (\emph{ensemble}) di molti \hl{alberi di decisione}. Ogni albero viene addestrato su una versione \key{leggermente diversa} dei dati e, a ogni domanda che pone, può guardare solo \key{alcune variabili scelte a caso}. Per fare una previsione si chiede a tutti gli alberi e si fa la \hl{media} (regressione) o si combinano le loro probabilità (classificazione).
\end{inbreve}

\subsection{La saggezza della folla}

\begin{intuizione}[La giuria dei barattoli di caramelle]
Chiedi a 500 persone quante caramelle ci sono in un barattolo. Ognuno sbaglia, qualcuno di molto. Eppure la \hl{media} delle risposte è spesso sorprendentemente vicina al valore vero. Funziona per un motivo preciso: gli errori sono \key{in parte casuali e in direzioni diverse}, quindi nella media si compensano.

C'è però una condizione: le persone non devono essersi messe d'accordo. Se tutti copiano il vicino, la media della folla è sbagliata quanto il primo che ha parlato. Il Random Forest è costruito proprio per avere \hl{tanti esperti diversi e il più possibile indipendenti}.
\end{intuizione}

Nel nostro caso gli «esperti» sono \hl{alberi di decisione}: modelli che fanno una sequenza di domande sì/no sulle variabili («il raggio del nucleo è maggiore di 16?», «l'età è minore di 40?») fino ad arrivare a una risposta. Un singolo albero ha due caratteristiche:
\begin{itemize}
  \item è \key{flessibile}: se lo lasci crescere fino in fondo, impara quasi perfettamente i dati di addestramento (basso \emph{bias});
  \item è \key{instabile}: basta cambiare pochi esempi e l'albero che ottieni può essere completamente diverso (alta \emph{varianza}).
\end{itemize}
Il Random Forest tiene il primo pregio e cura il secondo difetto: \hl{fa la media di molti alberi instabili}.

\subsection{Due ingredienti di casualità}

Per rendere gli alberi diversi tra loro, il Random Forest introduce due fonti di casualità, descritte da Breiman nel 2001 e usate anche in scikit-learn:

\begin{passi}
  \item \hl{Bootstrap dei dati} (\emph{bagging}). Ogni albero è addestrato su un campione estratto \key{con reinserimento} dal dataset: alcuni esempi compaiono due o tre volte, circa un terzo non compare affatto.
  \item \hl{Sottoinsieme casuale di feature a ogni nodo}. Quando un albero deve scegliere la domanda migliore, non può guardare tutte le variabili ma solo $m$ di esse, estratte a caso \key{ogni volta}. Così anche una variabile molto forte non domina tutti gli alberi.
\end{passi}

\begin{center}
\begin{tikzpicture}[x=1cm,y=1cm,
  dato/.style={draw=navy,fill=navy!8,rounded corners=3pt,minimum width=2.4cm,minimum height=.8cm,font=\headfont\bfseries\small,text=navy},
  boot/.style={draw=ambra,fill=ambra!12,rounded corners=3pt,minimum width=1.7cm,minimum height=.6cm,font=\headfont\footnotesize,text=ambra!50!black},
  alb/.style={draw=teal,fill=teal!10,rounded corners=3pt,minimum width=1.7cm,minimum height=.6cm,font=\headfont\footnotesize,text=teal!50!black},
  frec/.style={-{Stealth[length=5pt]},line width=.8pt,color=grigio}]
  \node[dato] (D) at (0,0) {Dataset};
  \foreach \i/\y in {1/1.6,2/.55,3/-.55,4/-1.6}{
    \node[boot] (b\i) at (3.3,\y) {bootstrap \i};
    \node[alb] (t\i) at (6.1,\y) {albero \i};
    \draw[frec] (D.east) -- (b\i.west);
    \draw[frec] (b\i) -- (t\i);}
  \node[font=\headfont\footnotesize,text=grigio] at (3.3,-2.3) {campioni diversi};
  \node[font=\headfont\footnotesize,text=grigio,align=center] at (6.1,-2.35) {$m$ feature a caso\\a ogni nodo};
  \node[draw=rosso,fill=rosso!10,rounded corners=3pt,minimum height=1.1cm,font=\headfont\bfseries\small,text=rosso!70!black,align=center] (V) at (9.4,0) {media /\\voto};
  \foreach \i in {1,...,4}{\draw[frec] (t\i.east) -- (V.west);}
  \node[font=\headfont\bfseries\small,text=navy] (P) at (11.6,0) {$\hat y$};
  \draw[frec] (V) -- (P);
\end{tikzpicture}
\end{center}
\didascalia{Schema generale: dallo stesso dataset si generano $B$ campioni bootstrap, si addestra un albero «randomizzato» su ciascuno e si combinano le risposte.}

\subsection{Perché la media aiuta: anticipo del risultato chiave}

Se $B$ alberi hanno ciascuno varianza $\sigma^2$ e tra due alberi qualunque la correlazione è $\rho$, la varianza della loro media è (lo dimostreremo nel capitolo 4)
\begin{formula}
$\displaystyle \Var\!\Big(\frac1B\sum_{b=1}^{B}T_b\Big)=\rho\,\sigma^2+\frac{1-\rho}{B}\,\sigma^2 .$
\end{formula}
Il secondo termine sparisce aggiungendo alberi; il primo no. Per questo \hl{non basta avere molti alberi: bisogna che siano poco correlati}. Il bootstrap e la scelta casuale delle feature servono esattamente a ridurre $\rho$.

\begin{mlbox}[Nel mondo reale]
Il Random Forest è tra i modelli più usati su \key{dati tabellari} (righe = esempi, colonne = variabili): diagnostica medica, rischio di credito, previsione della domanda, telerilevamento, bioinformatica. Funziona bene «così com'è», con pochi parametri da regolare, e per questo è spesso il primo modello robusto da provare dopo una regressione lineare o logistica.
\end{mlbox}

\begin{appunti}{Riepilogo del capitolo}
\begin{itemize}
  \item Random Forest = \hl{molti alberi di decisione} + \hl{media} delle risposte.
  \item Gli alberi singoli hanno basso bias ma alta varianza; la media riduce la varianza.
  \item La media funziona solo se gli alberi sono \key{diversi}: bootstrap dei dati + $m$ feature casuali a ogni nodo.
  \item Varianza della media: $\rho\sigma^2+(1-\rho)\sigma^2/B$ $\Rightarrow$ conta ridurre la correlazione $\rho$.
\end{itemize}
\end{appunti}
